% !TEX program = xelatex
% ============================================================
%  医学课程笔记整合 — 前言区（样式规范）
%  字体：Latin = Latin Modern（LaTeX 默认）；CJK = ctex 默认(fontset=mac, 宋/黑/楷)
%  主色：医学期刊式深墨绿蓝 + 青绿；少量暖色用于答案与装饰强调
% ============================================================

\usepackage{geometry}
\geometry{a4paper, top=2.0cm, bottom=1.9cm, left=1.6cm, right=1.6cm, headsep=0.4cm}
\setlength{\columnsep}{0.7cm}                 % 双栏栏间距
\setlength{\columnseprule}{0pt}               % 跨栏小节标题下避免竖线穿过标题区
\setlength{\headheight}{20pt}                 % fancyhdr 页眉高度，避免反复警告
\AtBeginDocument{\small}                      % 正文整体略缩小；标题/图表另行设定字号
\linespread{1.06}                             % 紧凑行距(长篇笔记)
% 窄栏断行容差(减少溢出)
\tolerance=3000
\setlength{\emergencystretch}{3.5em}
\hbadness=10000
\xeCJKsetup{CheckSingle=true}                 % 避免行首孤字

\usepackage{xcolor}
% 配色方案 D：医学期刊式编辑配色。深墨绿蓝建立层级，青绿保留医学/教学感，陶红铜色点到为止。
\definecolor{headchap}{HTML}{244B5A}       % 章标题：深墨绿蓝，接近印刷标题墨色
\definecolor{headteal}{HTML}{238C8F}       % 节/主色：低饱和医学青绿
\definecolor{headtealD}{HTML}{5D8388}      % 子标题/页眉：灰化青绿，降低压迫感
\definecolor{accentcopper}{HTML}{B36A4B}   % 小面积暖色：装饰点、提示边线
\definecolor{ruleteal}{HTML}{D7E5E7}       % 分隔/边框：浅蓝灰
\definecolor{bulletblue}{HTML}{6CA1A5}     % 列表符号：柔和青绿
\definecolor{supplbg}{HTML}{F8FBFA}        % 知识选读底色：近白微青
\definecolor{ansred}{HTML}{9B443C}         % 题库答案色：医学期刊常见的克制砖红
\definecolor{headtint}{HTML}{EEF6F6}       % 表格表头底：浅青灰
\definecolor{rowtint}{HTML}{F7FBFB}        % 表格隔行浅底

\usepackage{amsmath,amssymb}
\usepackage{textcomp}
\usepackage{gensymb}

% ---------- Unicode 符号统一映射（源文 OCR 残留的特殊符号）----------
% 算符/关系/箭头 → 数学模式；上下标 → \textsub/superscript；℃/° → 排版命令
\usepackage{newunicodechar}
% 运算符/箭头后允许断行(窄栏长公式折行,防溢出)
\newunicodechar{→}{\ensuremath{\rightarrow}\allowbreak}
\newunicodechar{←}{\ensuremath{\leftarrow}\allowbreak}
\newunicodechar{↑}{\ensuremath{\uparrow}\allowbreak}
\newunicodechar{↓}{\ensuremath{\downarrow}\allowbreak}
\newunicodechar{↔}{\ensuremath{\leftrightarrow}\allowbreak}
\newunicodechar{×}{\ensuremath{\times}\allowbreak}
\newunicodechar{÷}{\ensuremath{\div}\allowbreak}
\newunicodechar{±}{\ensuremath{\pm}\allowbreak}
\newunicodechar{∓}{\ensuremath{\mp}\allowbreak}
\newunicodechar{≤}{\ensuremath{\leq}\allowbreak}
\newunicodechar{≥}{\ensuremath{\geq}\allowbreak}
\newunicodechar{≠}{\ensuremath{\neq}\allowbreak}
\newunicodechar{≈}{\ensuremath{\approx}\allowbreak}
\newunicodechar{∝}{\ensuremath{\propto}}
\newunicodechar{∞}{\ensuremath{\infty}}
\newunicodechar{√}{\ensuremath{\surd}}
\newunicodechar{∑}{\ensuremath{\sum}}
\newunicodechar{∂}{\ensuremath{\partial}}
\newunicodechar{∶}{\ensuremath{:}}
\newunicodechar{∕}{/}
\newunicodechar{−}{\ensuremath{-}}
% 希腊字母
\newunicodechar{μ}{\ensuremath{\mu}}
\newunicodechar{α}{\ensuremath{\alpha}}
\newunicodechar{β}{\ensuremath{\beta}}
\newunicodechar{γ}{\ensuremath{\gamma}}
\newunicodechar{δ}{\ensuremath{\delta}}
\newunicodechar{ε}{\ensuremath{\varepsilon}}
\newunicodechar{η}{\ensuremath{\eta}}
\newunicodechar{θ}{\ensuremath{\theta}}
\newunicodechar{κ}{\ensuremath{\kappa}}
\newunicodechar{λ}{\ensuremath{\lambda}}
\newunicodechar{ν}{\ensuremath{\nu}}
\newunicodechar{π}{\ensuremath{\pi}}
\newunicodechar{ρ}{\ensuremath{\rho}}
\newunicodechar{σ}{\ensuremath{\sigma}}
\newunicodechar{τ}{\ensuremath{\tau}}
\newunicodechar{φ}{\ensuremath{\varphi}}
\newunicodechar{χ}{\ensuremath{\chi}}
\newunicodechar{ψ}{\ensuremath{\psi}}
\newunicodechar{ω}{\ensuremath{\omega}}
\newunicodechar{Δ}{\ensuremath{\Delta}}
\newunicodechar{Σ}{\ensuremath{\Sigma}}
\newunicodechar{Ω}{\ensuremath{\Omega}}
% 上下标
\newunicodechar{₀}{\textsubscript{0}}
\newunicodechar{₁}{\textsubscript{1}}
\newunicodechar{₂}{\textsubscript{2}}
\newunicodechar{₃}{\textsubscript{3}}
\newunicodechar{₄}{\textsubscript{4}}
\newunicodechar{₅}{\textsubscript{5}}
\newunicodechar{⁰}{\textsuperscript{0}}
\newunicodechar{¹}{\textsuperscript{1}}
\newunicodechar{²}{\textsuperscript{2}}
\newunicodechar{³}{\textsuperscript{3}}
\newunicodechar{⁻}{\textsuperscript{\textminus}}
\newunicodechar{⁺}{\textsuperscript{+}}
% 度/摄氏度/其他
\newunicodechar{℃}{\textdegree C}
\newunicodechar{°}{\textdegree}
\newunicodechar{‰}{\textperthousand}
\newunicodechar{•}{\textbullet}
\newunicodechar{·}{\ensuremath{\cdot}}
\newunicodechar{…}{\dots}
% 临床卷新增符号
\newunicodechar{∴}{\ensuremath{\therefore}}
\newunicodechar{∵}{\ensuremath{\because}}
\newunicodechar{≦}{\ensuremath{\leqq}}
\newunicodechar{≧}{\ensuremath{\geqq}}
\newunicodechar{●}{\textbullet}
\newunicodechar{➝}{\ensuremath{\rightarrow}}
\newunicodechar{﹥}{\textgreater{}}
\newunicodechar{⁹}{\textsuperscript{9}}
\newunicodechar{🈲}{（禁）}
% 罗马数字（分级用，源文 OCR 为 Unicode 罗马数字）
\newunicodechar{Ⅰ}{I}\newunicodechar{Ⅱ}{II}\newunicodechar{Ⅲ}{III}\newunicodechar{Ⅳ}{IV}
\newunicodechar{Ⅴ}{V}\newunicodechar{Ⅵ}{VI}\newunicodechar{Ⅶ}{VII}\newunicodechar{Ⅷ}{VIII}
\newunicodechar{Ⅸ}{IX}\newunicodechar{Ⅹ}{X}\newunicodechar{Ⅺ}{XI}\newunicodechar{Ⅻ}{XII}
% 注：U+2009 THIN SPACE / U+2002 EN SPACE 在转换管线中替换为普通空格

% pandoc 紧凑列表命令（裸 body 需自定义）
\providecommand{\tightlist}{\setlength{\itemsep}{0pt}\setlength{\parskip}{0pt}}

% 带圈数字 ①②③、性别符 ♂♀ 等路由到 CJK 字体（默认归 Latin 字体而缺字形）
\xeCJKDeclareCharClass{CJK}{"2460->"24FF, "2600->"26FF, "2160->"216F}

\usepackage{array}
\usepackage{booktabs}
\usepackage{makecell}
\usepackage{multirow}
\usepackage{longtable}
\usepackage{enumitem}
\usepackage{tabularray}
\usepackage{graphicx}
\usepackage{ragged2e}

\usepackage{tikz}
\usetikzlibrary{positioning,arrows.meta,shapes.geometric,shapes.misc,calc,fit,backgrounds,decorations.pathreplacing,patterns}
\usepackage{pgfplots}
\pgfplotsset{compat=1.18}

\usepackage[most]{tcolorbox}
\usepackage{cuted}
\usepackage{needspace}
\setlength{\stripsep}{11pt plus 2pt minus 1pt}

% ---------- 列表全局风格（紧凑、中文友好）----------
% 窄栏紧凑列表
\setlist[enumerate]{leftmargin=1.6em, itemsep=1pt, parsep=0.5pt, topsep=2pt}
\setlist[itemize]{leftmargin=1.4em, itemsep=1pt, parsep=0.5pt, topsep=2pt, label=\textcolor{bulletblue}{\textbullet}}
\setlist[enumerate,2]{label=(\arabic*), leftmargin=1.5em, itemsep=0.5pt}
\setlist[itemize,2]{leftmargin=1.3em, label=\textcolor{ruleteal}{\textendash}}
\setlist[itemize,3]{leftmargin=1.2em, label=\textcolor{ruleteal}{$\cdot$}}
\setlist[enumerate,3]{label=\alph*., leftmargin=1.4em}

% ---------- 章节标题样式（ctex 原生，避免 titlesec 冲突）----------
\setcounter{secnumdepth}{0}   % 仅章编号；节标题保留源文「第N节 ｜」文字
\setcounter{tocdepth}{1}      % 目录含章、节

\ctexset{
  chapter = {
    name        = {第,章},
    number      = {\chinese{chapter}},
    format      = {\sffamily\bfseries\color{headchap}},
    nameformat  = {\Huge},
    titleformat = {\Huge},
    aftername   = {\hspace{0.72em}},
    beforeskip  = {6pt},
    afterskip   = {26pt},
    pagestyle   = {plain},
  },
  section = {
    format      = {\sffamily\bfseries\Large\color{headteal}},
    beforeskip  = {2.2ex},
    afterskip   = {1.2ex},
  },
  subsection = {
    format      = {\sffamily\bfseries\large\color{headtealD}},
    beforeskip  = {1.6ex},
    afterskip   = {0.8ex},
  },
}

% ---------- 课程小节：标题跨双栏居中，标题下正文仍为双栏 ----------
% 只由生成脚本用于课程笔记正文；题库附录保留普通 \section。
\newcommand{\courseornament}{%
  \begin{tikzpicture}[baseline=-0.45ex]
    \draw[ruleteal, line width=0.45pt, line cap=round] (-2.85,0) -- (-0.42,0);
    \draw[ruleteal, line width=0.45pt, line cap=round] (0.42,0) -- (2.85,0);
    \fill[ruleteal] (-3.05,0) circle[radius=0.65pt];
    \fill[ruleteal] (3.05,0) circle[radius=0.65pt];
    \draw[accentcopper, line width=0.55pt] (0,0.095) -- (0.095,0) -- (0,-0.095) -- (-0.095,0) -- cycle;
    \fill[accentcopper] (0,0) circle[radius=0.8pt];
  \end{tikzpicture}%
}

\newcommand{\coursesection}[1]{%
  \if\relax\detokenize{#1}\relax\else
  \par\addvspace{2.4ex}%
  \Needspace{7\baselineskip}%
  \begingroup
    \phantomsection
    \addcontentsline{toc}{section}{#1}%
    {\centering\sffamily\bfseries\Large\color{headteal} #1\par}%
    \vspace{0.35ex}%
    {\centering\courseornament\par}%
  \endgroup
  \par\nobreak\addvspace{0.9ex}%
  \fi
}

\newcommand{\coursesubsection}[1]{%
  \if\relax\detokenize{#1}\relax\else
  \par\addvspace{1.5ex}%
  \Needspace{4\baselineskip}%
  {\noindent\sffamily\bfseries\large\color{headtealD} #1\par}%
  \par\nobreak\addvspace{0.45ex}%
  \fi
}

\newcommand{\coursesubsubsection}[1]{%
  \if\relax\detokenize{#1}\relax\else
  \par\addvspace{1.1ex}%
  \Needspace{3\baselineskip}%
  {\noindent\sffamily\bfseries\normalsize\color{headtealD} #1\par}%
  \par\nobreak\addvspace{0.25ex}%
  \fi
}

\newcommand{\courseparagraph}[1]{%
  \if\relax\detokenize{#1}\relax\else
  \par\addvspace{0.8ex}%
  \Needspace{2\baselineskip}%
  {\noindent\sffamily\bfseries\small\color{accentcopper} #1\par}%
  \par\nobreak\addvspace{0.15ex}%
  \fi
}

% 疼痛诊疗学的模块层标题：用于“诊断学/治疗学总论/治疗学分论”等中间层。
\newcommand{\coursegroup}[1]{%
  \if\relax\detokenize{#1}\relax\else
  \par\addvspace{2.0ex}%
  \Needspace{6\baselineskip}%
  \begingroup
    \phantomsection
    \addcontentsline{toc}{section}{#1}%
    {\centering\sffamily\bfseries\large\color{headtealD} #1\par}%
    \vspace{0.35ex}%
    {\centering
      \begin{tikzpicture}[baseline=-0.45ex]
        \draw[ruleteal, line width=0.36pt, line cap=round] (-2.1,0) -- (-0.28,0);
        \draw[ruleteal, line width=0.36pt, line cap=round] (0.28,0) -- (2.1,0);
        \draw[accentcopper, line width=0.45pt] (0,0.075) -- (0.075,0) -- (0,-0.075) -- (-0.075,0) -- cycle;
      \end{tikzpicture}\par}%
  \endgroup
  \par\nobreak\addvspace{0.6ex}%
  \fi
}

\providecommand{\painsectiontoc}[1]{#1}

% 疼痛诊疗学的具体节标题：作为正文局部标题显示，不再创建全书私有目录层。
\newcommand{\painsection}[1]{\coursesubsection{#1}}

% 章下方一条浅青横线（呼应参考 PDF 标题区）
\let\oldchapterheadendhook\relax

% ---------- 目录：点线引导 + 章节着色 ----------
\usepackage{titletoc}
% 章：青色、加粗、带点线引导至页码
\titlecontents{chapter}[0pt]
  {\addvspace{6pt}\sffamily\bfseries\color{headchap}}
  {}{}
  {\hspace{0.45em}\color{ruleteal}\titlerule*[0.78pc]{.}\contentspage}
% 节：缩进、常规色、点线
\titlecontents{section}[2.2em]
  {\addvspace{1pt}\small}
  {}{}
  {\hspace{0.45em}\color{ruleteal}\titlerule*[0.78pc]{.}\normalcolor\contentspage}
% 兼容旧 .toc 中可能残留的 painsection 条目；新构建不再写入该目录层。
\titlecontents{painsection}[4.6em]
  {\addvspace{0pt}\footnotesize\color{black!70}}
  {}{}
  {\hspace{0.4em}\color{ruleteal}\titlerule*[0.82pc]{.}\normalcolor\contentspage}
\makeatletter
\providecommand*{\toclevel@painsection}{2}
\@ifundefined{ttll@painsection}{\gdef\ttll@painsection{1}}{}
\makeatother
% 部分/附录分隔页：居中大字
\titlecontents{part}[0pt]
  {\addvspace{12pt}\centering\sffamily\bfseries\large\color{headchap}}
  {}{}{}

% ---------- 页眉页脚 ----------
\usepackage{fancyhdr}
\pagestyle{fancy}
\fancyhf{}
\fancyhead[C]{\small\sffamily\color{headtealD}\leftmark}
\fancyfoot[C]{\small\color{headtealD}\thepage}
\renewcommand{\headrulewidth}{0.4pt}
\renewcommand{\footrulewidth}{0pt}
\renewcommand{\headrule}{\hbox to\headwidth{\color{ruleteal}\leaders\hrule height \headrulewidth\hfill}}

% ---------- 分隔页（无编号·青色居中；双栏文档下单栏整页）----------
\newcommand{\volpart}[1]{%
  \clearpage \onecolumn
  \phantomsection
  \addcontentsline{toc}{part}{#1}%
  \markboth{#1}{#1}%
  \thispagestyle{empty}%
  \null\vfill
  {\centering\sffamily\bfseries\fontsize{32}{38}\selectfont\color{headchap} #1\par}%
  \vskip 1.2em
  {\centering\color{accentcopper}\rule{0.36\textwidth}{0.9pt}\par}%
  \vfill\clearpage \twocolumn
}

% ---------- 普通课程章：保留 ctex 章样式，同时修正 fancyhdr 页眉 ----------
\newcommand{\coursechapter}[1]{%
  \chapter{#1}%
  \markboth{#1}{#1}%
}

% ---------- 可选整页章入口；默认正文使用 \coursechapter，保持全书一致 ----------
\newcommand{\disciplinechapter}[1]{%
  \clearpage \onecolumn
  \refstepcounter{chapter}%
  \phantomsection
  \addcontentsline{toc}{chapter}{\protect\numberline{\thechapter}#1}%
  \markboth{#1}{#1}%
  \thispagestyle{empty}%
  \null\vfill
  {\centering\sffamily\bfseries\Large\color{headchap} 第\chinese{chapter}章\par}%
  \vskip 1.0em
  {\centering\sffamily\bfseries\fontsize{36}{44}\selectfont\color{headchap} #1\par}%
  \vskip 1.2em
  {\centering\color{accentcopper}\rule{0.42\textwidth}{0.9pt}\par}%
  \vfill\clearpage \twocolumn
}

% ---------- 名词解释 / 问答题 / 难点集萃：带色标签块 ----------
% 用法：\noteblock{名词解释} 后接列表内容
\newcommand{\noteblock}[1]{%
  \par\addvspace{1.2ex}%
  \noindent\begin{tikzpicture}[baseline]
    \node[anchor=west, inner sep=0pt] (a) at (0,0)
      {\color{accentcopper}\rule{3.2pt}{1.05em}};
    \node[anchor=west, inner sep=0pt, xshift=6pt] at (a.east)
      {\sffamily\bfseries\large\color{headtealD} #1};
  \end{tikzpicture}\par\nobreak\addvspace{0.5ex}%
}

% ---------- 补充内容 · 知识选读：浅青旁注框 ----------
\newtcolorbox{supplbox}[1][补充内容 · 知识选读]{%
  enhanced, breakable,
  colback=supplbg, colframe=ruleteal, boxrule=0.4pt, arc=1.2mm,
  left=2.6mm, right=2.6mm, top=1.6mm, bottom=1.6mm,
  borderline west={1.8pt}{0pt}{accentcopper},
  fonttitle=\sffamily\bfseries\small\color{white},
  coltitle=white, colbacktitle=headtealD,
  attach boxed title to top left={xshift=2mm, yshift=-2mm},
  boxed title style={arc=0.6mm, boxrule=0pt},
  title={#1},
}

% ---------- 题库：题目 / 答案 ----------
% \qitem{题号}{题干} ... \begin{choices}...\end{choices} ... \ans{X}
\newcommand{\qitem}[2]{\par\addvspace{0.6ex}\noindent\textbf{#1.}~#2\par\nobreak}
\newcommand{\ans}[1]{\par\nopagebreak{\sffamily\bfseries\color{ansred}答案：#1}\par\addvspace{0.4ex}}
\newcommand{\qreview}[1]{\par\nopagebreak\noindent{\footnotesize\sffamily\color{headtealD}\textbf{\textcolor{ansred}{疑点：}}#1}\par\addvspace{0.4ex}}
\newlist{choices}{enumerate}{1}
\setlist[choices]{label=\Alph*、, leftmargin=1.9em, itemsep=0.3pt, topsep=1pt, parsep=0pt}

% ---------- 表格通用样式（浅青表头 + booktabs 横线）----------
\UseTblrLibrary{booktabs}
\NewTblrEnviron{mtbl}
\SetTblrInner[mtbl]{rowsep=2.2pt, colsep=5pt}

% ---------- 节间分隔符（替换源文 --- ）----------
\newcommand{\coursediv}{\par\addvspace{0.6ex}\centerline{\color{ruleteal}\rule{0.30\textwidth}{0.6pt}}\par\addvspace{0.6ex}}
\let\anesdiv\coursediv

% ---------- 引用块（定义性内容，浅青左标）----------
\renewenvironment{quote}
  {\begin{list}{}{\setlength{\leftmargin}{1.4em}\setlength{\rightmargin}{0em}}\item[]%
   \begingroup\color{black}}
  {\endgroup\end{list}}

% ---------- 神经肌肉监测柱状助手 ----------
\newcommand{\nmbar}[1]{% #1=逗号分隔高度(0为空隙)
  \begin{tikzpicture}[baseline=-0.5pt,x=1.15mm,y=2.7mm]
    \foreach \h [count=\i from 0] in {#1}{%
      \pgfmathparse{\h>0 ? 1 : 0}%
      \ifnum\pgfmathresult=1 \fill[headteal] (\i,0) rectangle ++(0.62,\h);\fi}
  \end{tikzpicture}}

\usepackage[hidelinks,bookmarksnumbered]{hyperref}
\usepackage{bookmark}
